\section{Terrain and Structures}

Terrain features and structures are important components of the battlefield. A base is considered to be within a terrain feature if the centre of the base is inside it.

\subsection{Effects on Movement}

A terrain feature affects the movement of bases according to its type and class.

\begin{description}
    \item[Impassable:] These terrain features prevent bases from moving through them. Bases must move around them or, if they possess an appropriate ability, fly over them. Cliffs and deep water are good examples of Impassable terrain.

    \item[Difficult Terrain:] These terrain features make movement more difficult. Bases move at half speed while moving through them. Ruins and forests are good examples of Difficult Terrain.

    \item[Dangerous Terrain (X+/AP Y):] These terrain features are hazardous to cross. A base that moves through Dangerous Terrain, or begins the Combat Phase within it, suffers a hit on a roll of X+ with AP~Y. This is treated as a shooting attack that ignores cover saves.

    If the AP value is listed as ``None'', no saving throw is permitted. A base makes no more than one roll per turn for each individual area of Dangerous Terrain. Minefields and swamps are good examples of Dangerous Terrain.
\end{description}

\subsection{Effects on Line of Sight}

Terrain features and structures affect line of sight. Each terrain feature and structure must be assigned a class for the purpose of determining line of sight.

If a terrain feature does not have a uniform height, or if it is particularly large, it may be divided into several separate features with different classes.

\begin{description}
    \item[Line-of-Sight Blocking:] These terrain features completely block line of sight. Rocks and hills are good examples.

    If bases can occupy a Line-of-Sight Blocking terrain feature, line of sight is blocked as soon as it passes through the terrain. A base must therefore be positioned at the edge of the terrain feature to have line of sight beyond it.

    \item[Obstructing:] These terrain features do not block line of sight but provide a cover save. If the line of sight passes through Obstructing terrain, the target receives its cover save even if the target base is not inside the terrain feature.

    If a base is inside an Obstructing terrain feature, that feature does not affect the base's own line of sight. Ruins are good examples of Obstructing terrain.
\end{description}

\subsection{Cover Saves}

Some terrain features provide a cover save. This is a fixed saving throw that is not modified by AP and is made before any normal saving throws.

Cover saves are cumulative. Each additional source of cover improves the best available cover save by 2.

For example, if a shooting attack passes through two ruins, the resulting cover save is 3+: the initial 5+ cover save is improved by 2.

\subsection{Terrain Size}
A terrain feature's size is represented by its class, which must be agreed upon before the battle begins.

As a guideline, a terrain feature may be assigned one class for every 1.5~cm of its actual height. Adjust this guideline to suit the scale of your miniatures.

\subsection{Effects on Assaults}

Some terrain features provide an AF bonus to bases occupying them. This bonus always applies, regardless of whether the opposing base also occupies the same terrain feature.

\subsection{Structures}

Structures are enclosed, easily defended terrain features, such as buildings and bunkers. They function differently from other types of terrain.

Bases that enter a structure are considered to be \textit{garrisoned}.

\subsubsection{Structure Size}

Some structures are large and should be divided into several smaller structures.

It is recommended that each structure be no larger than approximately 10~cm by 10~cm, with a capacity of between 6 and 12 bases.

\subsubsection{Entering and Leaving a Structure}

Structures have a limited capacity and can therefore contain only a limited number of bases.

A base may enter a structure only if there is sufficient space available, there are no enemy bases inside it, and no assault is currently being fought against its garrison.

Entering or leaving a structure costs 5~cm of movement. A base cannot both enter and leave the same structure during the same turn.

Bases may enter or leave through any edge of the structure.

\subsubsection{Line of Sight}

For line-of-sight purposes, bases inside a structure are treated as having the same class as the structure.

The exact positions of bases inside the structure are irrelevant. Line of sight may be drawn from any point on the structure. This applies both to garrisoned bases drawing line of sight and to enemy bases targeting them.

\subsubsection{Zone of Control}

A garrison's zone of control extends around the entire perimeter of the structure, as though the centre of each garrisoned base were positioned at its edge. This applies regardless of the number of bases inside the structure.

\subsubsection{Assault}

The attacker engages bases within the structure according to the normal rules. Each attacking base engages one defending base, and the defending player chooses which defending base is engaged by each attacker.

The defending player may then choose to Consolidate with any unengaged bases, while remaining inside the structure if desired.

When resolving an assault, garrisoned bases are not pushed back by a Regroup move if no enemy base is able to enter the structure.

If enemy bases end their movement close enough to trigger a Consolidation, garrisoned bases may begin their Consolidation move from any point on the structure. However, bases that Consolidate must leave the structure and are no longer considered garrisoned.

\subsubsection{Destroying Structures}

Only weapons with the \textit{Damages Buildings} ability may target and damage structures. Shooting attacks targeting a structure receive a $+2$ bonus to their To-Hit rolls.

The structure makes its saving throws using 2D6, modified by the AP value specified by the firing weapon's \textit{Damages Buildings} ability.

When a structure is reduced to zero Wounds, it is destroyed and replaced with a ruin of the same size. Any bases inside it may make an Emergency Exit roll, as described by the \textit{Transport} ability, modified by the total amount of damage inflicted by the activation that destroyed the structure.

If the firing detachment has several attacks capable of damaging buildings, add together all the damage inflicted by that detachment. If the bases survive their Emergency Exit rolls, their owner places them anywhere within the resulting ruin.

If a template weapon destroys the structure, first resolve its hits against the garrisoned bases and then make the structure's saving throw.

If the firing detachment has attacks other than the attack capable of damaging the structure, resolve the standard shooting attacks first and the building-damaging attack afterwards.

\subsection{Purchasing Fortifications}

Some armies may purchase fortifications. These are included in the army list and deployed during the battle preparation sequence.

The points cost of any purchased fortifications is automatically awarded to the opposing player at the beginning of the battle.

\subsection{Terrain Types}

\begin{description}
    \item[Barricades:] Barricades are obstacles that slow movement and provide light protection. A Class 1 base in contact with a barricade receives a 5+ cover save and a $+1$ AF bonus.

    Class 1--3 bases crossing a barricade lose 5~cm of movement.

    \item[Hills:] Hills do not impede movement and affect only line of sight. They may block line of sight, and their class is added to the class of any base positioned on them when determining line of sight.

    \item[Forests:] A forest is a Line-of-Sight Blocking terrain feature that provides a 5+ cover save to every base occupying it whose class is equal to or lower than the forest's designated class.

    If a line of sight passes through a forest, it is blocked. If a base is inside a forest, a valid line of sight must pass directly from the base without passing through any other part of the forest.

    Overwatch Fire is always possible immediately before an enemy base makes contact. If two bases within the same forest wish to target one another, their line-of-sight distance is limited to 12~cm.

    \item[Trenches:] Trenches provide a 4+ cover save and a $+2$ AF bonus. Only infantry bases may enter them.

    Class 1--3 bases crossing a trench lose 5~cm of movement.

    \item[Swamps:] Swamps provide a 5+ cover save to Class 1 and 2 bases. They count as Difficult Terrain for all classes except Titans and Praetorians.

    \item[Rivers:] Rivers count as Difficult Terrain and \textit{Dangerous Terrain (5+/None)} for infantry. Apart from Titans and Praetorians, no other bases may cross them.

    \item[Ruins:] Ruins are Obstructing terrain features that provide a 5+ cover save, generally to Class 1--4 bases. The applicable class must be determined according to the nature of the ruin.

    Ruins also count as Difficult Terrain for Cavalry, Vehicles, Super-heavy Vehicles, and Knights.

    \item[Craters:] Craters provide a 5+ cover save, generally to Class 1--3 bases. The applicable class must be determined according to the nature of the crater.

    Craters also count as Difficult Terrain for Cavalry, Vehicles, Super-heavy Vehicles, and Knights.

    \item[Cliffs:] Cliffs are Impassable unless the moving base has at least twice the class of the cliff.

    \item[Structures:] Bases inside a structure benefit from the garrison rules. Structures provide a cover save and an AF bonus according to their type.

    \item[Light Structures:] Light structures provide a 5+ cover save and a $+1$ AF bonus. They have a 5+ Save and 2 Wounds.

    \item[Standard Structures:] Standard structures provide a 4+ cover save and a $+2$ AF bonus. They have a 4+ Save and 3 Wounds.

    \item[Bastion:] Bastions are large bunkers with a Titan-class weapon mounted on the roof. This weapon is included in their points cost and gains the \textit{Turret} and \textit{Autonomous Defences} abilities.

    A Bastion provides a 4+ cover save and a $+2$ AF bonus. It has a 3+ Save, 3 Wounds, and a transport capacity of 10 infantry bases.

    \item[Bunker:] Bunkers are either large bunkers or complexes consisting of several smaller defensive positions.

    A Bunker provides a 4+ cover save and a $+2$ AF bonus. It has a 3+ Save, 3 Wounds, and a transport capacity of 10 infantry bases.

    \item[Artillery Emplacement:] Artillery Emplacements are large bunkers designed to protect vehicles. They provide a 4+ cover save and a $+2$ AF bonus, and have a 3+ Save and 3 Wounds.

    Infantry, Walker, and Cavalry bases may enter them, but receive only the 4+ cover save.

    An Artillery Emplacement has a capacity of either three Vehicle bases or six Infantry, Walker, or Cavalry bases.

    \item[Minefields:] A minefield is \textit{Dangerous Terrain (4+/AP 2)} that ignores Shields and Holo-fields for every base not at altitude.

    A base is affected as soon as any part of it comes into contact with the minefield, even if the centre of the base has not entered it.

    A base with multiple Wounds suffers one attack for each of its initial Wounds. These attacks always hit Location~1 on its hit-location chart, representing its lowest section.
\end{description}

\begin{table}[ht]
    \centering

    \resizebox{\textwidth}{!}{%
        \begin{tblr}{
            colspec={
                Q[c,m] Q[c,m] Q[c,m] Q[c,m] Q[c,m] Q[c,m]
                Q[c,m] Q[c,m] Q[c,m] Q[c,m] Q[c,m]
            },
            cells={
                halign=c,
                valign=m
            },
            hlines,
            vlines,
            row{1}={
                bg=black,
                fg=white,
                font=\bfseries,
                ht=12mm,
                valign=m
            },
            row{3,5,7,9,11}={bg=gray!20}
        }
            Terrain &
            Infantry &
            Walkers &
            Cavalry &
            Vehicles &
            \SetCell{c,m}{Knights and\\Super-heavies} &
            \SetCell{c,m}{Titans and\\Praetorians} &
            Line of Sight &
            \SetCell{c,m}{Save\\(2D6)} &
            \SetCell{c,m}{Cover\\Save} &
            AF Bonus
            \\

            \textbf{Forest} &
            \SetCell[c=2]{c,m} Normal & &
            \SetCell[c=3]{c,m} Difficult & & &
            Normal &
            \SetCell{c,m}{Blocking\\Class 4} &
            Indestructible &
            5+ &
            0
            \\

            \textbf{Swamp} &
            \SetCell[c=2]{c,m} Difficult & &
            \SetCell[c=2]{c,m}{Difficult,\\Dangerous (6+/None)} & &
            Difficult &
            Normal &
            No Effect &
            Indestructible &
            \SetCell{c,m}{5+\\Classes 1 and 2} &
            0
            \\

            \textbf{River and Lake} &
            \SetCell{c,m}{Difficult,\\Dangerous (5+/None)} &
            \SetCell[c=4]{c,m} Impassable & & & &
            Normal &
            No Effect &
            Indestructible &
            -- &
            0
            \\

            \SetCell{c,m}{\textbf{Ruins and}\\\textbf{Craters}} &
            \SetCell[c=2]{c,m} Normal & &
            \SetCell[c=3]{c,m} Difficult & & &
            Normal &
            \SetCell{c,m}{Obstructing\\Classes 1--7} &
            Indestructible &
            5+ &
            0
            \\

            \textbf{Barricade} &
            \SetCell[c=4]{c,m}
                Costs 5~cm of movement to cross & & & &
            \SetCell[c=2]{c,m} Normal & &
            \SetCell{c,m}{Blocking\\Class 1} &
            \SetCell{c,m}{7+\\1 Wound per 10~cm} &
            5+ &
            $+1$
            \\

            \textbf{Trench} &
            \SetCell[c=4]{c,m}
                Costs 5~cm of movement to cross & & & &
            \SetCell[c=2]{c,m} Normal & &
            No Effect &
            -- &
            4+ &
            $+2$
            \\

            \SetCell{c,m}{\textbf{Light}\\\textbf{Structure}} &
            Normal &
            \SetCell[c=5]{c,m} Impassable & & & & &
            \SetCell{c,m}{Blocking\\Classes 4--7} &
            \SetCell{c,m}{5+\\2 Wounds} &
            5+ &
            $+1$
            \\

            \SetCell{c,m}{\textbf{Standard}\\\textbf{Structure}} &
            Normal &
            \SetCell[c=5]{c,m} Impassable & & & & &
            \SetCell{c,m}{Blocking\\Classes 4--7} &
            \SetCell{c,m}{4+\\3 Wounds} &
            4+ &
            $+2$
            \\

            \textbf{Bunker} &
            Normal &
            \SetCell[c=5]{c,m} Impassable & & & & &
            \SetCell{c,m}{Blocking\\Class 3} &
            \SetCell{c,m}{3+\\3 Wounds} &
            3+ &
            $+2$
            \\

            \SetCell{c,m}{\textbf{Artillery}\\\textbf{Emplacement}} &
            \SetCell[c=4]{c,m}
                May enter; otherwise Impassable & & & &
            \SetCell[c=2]{c,m} Impassable & &
            \SetCell{c,m}{Blocking\\Class 3} &
            \SetCell{c,m}{3+\\3 Wounds} &
            4+ &
            $+2$
            \\

            \textbf{Bastion} &
            Normal &
            \SetCell[c=5]{c,m} Impassable & & & & &
            \SetCell{c,m}{Blocking\\Class 3} &
            \SetCell{c,m}{3+\\4 Wounds} &
            3+ &
            $+2$
        \end{tblr}%
    }
\end{table}